% ---
% id: statistics/rem:two-tags
% kind: rem
% title: "The two tags and the consumer's certified geometries"
% summary: >-
%   The two tags are exactly the two geometries over which the consuming engine
%   certifies the minimum of a linear functional, and the canonical box form is
%   a sub-form of the polytope, not a third tag. A consumer that ignores the
%   recorded bounds still reads a well-formed polytope.
% prominence: standard
% status: accepted
% uses: []
% mentions: [def:taxonomy, thm:box-oracle, req:box-bounds-key, meat-core/frank-wolfe/rem:acc-side]
% assumes: []
% refs: []
% provenance: newly-derived
% lean:
% ---

\begin{remark}\label{rem:two-tags}
The two tags of Definition~\ref{def:taxonomy} correspond one to one to the two
instantiations of \xref{meat-core/frank-wolfe/rem:acc-side}, the geometries
for which that remark certifies the acceptance-side minimum, the minimum of a
linear functional over the acceptance set: (a) singleton, (b)
polytope $\{\mathbf{q}\ge0,\ \sum_{\hat c}q(\hat c)=1,\
B\mathbf{q}\le\mathbf{b}\}$ for exact $B$ and $\mathbf{b}$. That remark
certifies no set of any other form. The canonical box form is a
\emph{sub-form} of $\mathrm{polytope}$, not a third tag: it adds a recorded
pair $(\qlo,\qhi)$ (Requirement~\ref{req:box-bounds-key}) that licenses the
closed-form oracle of Theorem~\ref{thm:box-oracle}, and a \tref{consumer} that
ignores the recorded pair still reads a well-formed $\mathrm{polytope}$.
\end{remark}
